% =====================================================================
% Chapter 3 -- Methodology Overview
% =====================================================================
\documentclass[../preamble.tex]{subfiles}
\begin{document}

\chapter{Methodology Overview}
\label{ch:method-overview}

\section{Design intent}
\label{sec:method-intent}

The \PEARL{} pipeline was conceived as a sequence of seven phases that build
on each other, where each phase produces an artefact that feeds the next.
The pipeline is config-driven: every experimental variation is a YAML
override. The phases are described in the chapters that follow; this chapter
presents the data-flow structure that ties them together.

\begin{figure}[htbp]
  \centering
  \begin{tikzpicture}[
      node distance=0.6cm and 0.8cm,
      phase/.style={draw, rectangle, minimum height=0.9cm, minimum width=2.5cm, align=center, fill=blue!8},
      arr/.style={-{Latex[length=2mm]}},
    ]
    \node[phase] (p1) {Phase 1\\Data ingest + dedup};
    \node[phase, right=of p1] (p2) {Phase 2\\Split (80/10/10)};
    \node[phase, right=of p2] (p3) {Phase 3\\Preprocessing\\(CLAHE + denoiser)};
    \node[phase, right=of p3] (p4) {Phase 4\\Augment + train\\foundation};
    \node[phase, right=of p4] (p5) {Phase 5\\HPO + retrain\\(top 3)};
    \node[phase, below=of p2] (p6) {Phase 6\\Ensemble + calibration};
    \node[phase, below=of p5] (p7) {Phase 7\\Uncertainty + XAI};
    \draw[arr] (p1) -- (p2);
    \draw[arr] (p2) -- (p3);
    \draw[arr] (p3) -- (p4);
    \draw[arr] (p4) -- (p5);
    \draw[arr] (p5) -- (p6);
    \draw[arr] (p5) -- (p7);
    \draw[arr, dashed] (p6) -- (p7);
  \end{tikzpicture}
  \caption{Data-flow overview of the seven pipeline phases. Phase 6 (calibration)
           and Phase 7 (uncertainty and explanation) operate on the same
           checkpoints produced by Phases 4--5.}
  \label{fig:pipeline-overview}
\end{figure}

\section{Three training stages}
\label{sec:method-stages}

The thesis distinguishes three training stages that appear repeatedly in the
later chapters:

\begin{enumerate}
  \item \textbf{Foundation stage.} Each of nine backbones is fine-tuned on the
        Figshare training split under a fixed preprocessing pipeline. The
        result is eighteen foundation checkpoints (nine architectures $\times$
        two denoisers). Foundation checkpoints are evaluated on the Figshare
        held-out test set and, zero-shot, on the full PCOSgen corpus.
  \item \textbf{Fine-tune stage.} A subset of the strongest foundation
        checkpoints is fine-tuned on the PCOSgen training pool, with the
        held-out PCOSgen test set reserved for external evaluation.
  \item \textbf{Ensemble stage.} The three highest-AUC fine-tuned checkpoints
        are combined through probability averaging and then recalibrated by
        two-pass temperature scaling.
\end{enumerate}

The three stages map onto the three-panel confusion-matrix figure shown in
\cref{fig:confusion-stages} in the Results chapter.

\section{Inputs and outputs per stage}
\label{sec:method-io}

Each stage consumes the previous stage's artefact and produces a new
artefact:

\begin{table}[htbp]
  \centering
  \small
  \caption{Inputs and outputs of each training stage.}
  \label{tab:method-io}
  \begin{tabular}{lll}
    \toprule
    \textbf{Stage} & \textbf{Inputs} & \textbf{Outputs} \\
    \midrule
    Foundation & Figshare split + preprocessing & 18 checkpoints + sweep matrix \\
    Fine-tune & Foundation checkpoints + PCOSgen train pool & 18 fine-tuned checkpoints + per-arch metrics \\
    Ensemble & Top-3 fine-tuned checkpoints & probability file, calibration artefacts \\
    \bottomrule
  \end{tabular}
\end{table}

\section{Role of the design versus executed experiment}
\label{sec:method-design-vs-run}

The v3 methodology specification described several features that are not
present in the executed final run. The most important deviations are:

\begin{itemize}
  \item \textbf{Backbone freezing.} The methodology spec fixed a default
        \texttt{freeze\_fraction=0.60}, with HPO over $\{0.40, 0.50, 0.60,
        0.70\}$. The executed top-3 fine-tunes all selected
        \texttt{freeze\_fraction=0} after Optuna search.
  \item \textbf{SWA, EMA, k-fold CV.} The methodology spec describes SWA,
        EMA, and per-architecture $k$-fold cross-validation. These were
        implemented in the trainer scaffold but were not used to produce the
        reported results, which are based on single-split fine-tunes.
  \item \textbf{XAI methods.} The proposal listed Grad-CAM, LRP, and SHAP.
        Only Grad-CAM was executed on the fine-tuned checkpoints.
\end{itemize}

These deviations are stated up front here so the reader can interpret the
remaining chapters accurately. The full intended-versus-executed table is in
\cref{app:intended-executed}.

\section{Reproducibility posture}
\label{sec:method-repro}

Every result reported in this thesis is reproducible from the artefacts in
the repository. The reproducibility recipe, including package versions,
hardware, wall-clock estimates, and command-line invocations, is collected
in \cref{app:reproducibility}.

\end{document}